\chapter{No Arbitrage and the Fundamental Theorems}\label{ch:dv:no-arbitrage-and-the-fundamental-theorems}

Four S\&P~500 index options sit in one order: a call bought and a call sold,
a put sold and a put bought, at the strikes 5\,000 and 6\,000, all expiring in
a year. Whatever the index does, the four together pay exactly 1\,000 index
points at expiry, USD~100\,000 per set of contracts. The package trades at
958.90 points. Nobody on the desk thinks of it as an option position: it is a
one-year zero-coupon bond, guaranteed by the clearing house, and 958.90 is an
interest rate of 4.20\%. Some funds lend cash this way, others borrow; the
exchange publishes the market's size. Behind the trade sits the idea on which
every chapter of this book rests: two portfolios with the same payoff must
have the same price, or someone takes the difference for nothing. This
chapter makes the idea precise, turns it into prices of states of the world,
and proves the two theorems that tell us when a price is forced and when it
is only bounded.

\section{Arbitrage and the law of one price}

Fix a date $0$, a date $T$, and a finite set of states $\omega_1,\dots,\omega_S$,
one of which will occur at $T$. There are $N$ traded assets. Asset $j$ costs
$p_j$ today and pays $D_{sj}$ in state $s$: the market is the price vector
$p\in\R^N$ and the $S\times N$ payoff matrix $D$. A portfolio
$\theta\in\R^N$ (positive entries bought, negative sold) costs $p\cdot\theta$
and pays the vector $D\theta$.

\begin{definition}[Contingent claim and replication]\label{def:dv:no-arbitrage-and-the-fundamental-theorems:claim}
A \emph{contingent claim}\index{contingent claim} is a payoff at $T$ that
depends on the state: a vector $g\in\R^S$ in the one-period model, a random
variable $g(S_T)$ or $g(S_t, t\le T)$ in general. A \emph{replicating
portfolio}\index{replicating portfolio} for $g$ is a portfolio of traded assets
whose payoff equals $g$ in every state: $D\theta = g$.
\end{definition}

\begin{definition}[Arbitrage]\label{def:dv:no-arbitrage-and-the-fundamental-theorems:arbitrage}
An \emph{arbitrage}\index{arbitrage} is a portfolio that costs nothing or
less, never pays a negative amount, and is not identically worthless: $p\cdot
\theta\le0$, $D\theta\ge0$, and at least one of the inequalities strict ($p\cdot
\theta<0$, or $D\theta\ne0$). The market has \emph{no arbitrage} when no such
portfolio exists.
\end{definition}

\begin{definition}[Law of one price]\label{def:dv:no-arbitrage-and-the-fundamental-theorems:loop}
The \emph{law of one price}\index{law of one price} holds when two portfolios
with the same payoff have the same cost: $D\theta=D\theta'$ implies $p\cdot
\theta=p\cdot\theta'$. Under it, every replicable claim has a unique price, the
cost of any of its \omterm{def:dv:no-arbitrage-and-the-fundamental-theorems:claim}{replicating portfolios}.
\end{definition}

No \omterm{def:dv:no-arbitrage-and-the-fundamental-theorems:arbitrage}{arbitrage} implies the \omterm{def:dv:no-arbitrage-and-the-fundamental-theorems:loop}{law of one price} (if two portfolios with equal payoffs
had different costs, buying the cheap one and selling the dear one would be an
\omterm{def:dv:no-arbitrage-and-the-fundamental-theorems:arbitrage}{arbitrage}), but not the converse: a market in which an asset paying $(1,0)$
costs zero satisfies the \omterm{def:dv:no-arbitrage-and-the-fundamental-theorems:loop}{law of one price} and still gives something for
nothing (\cref{exo:dv:no-arbitrage-and-the-fundamental-theorems:4}). \omterm{def:dv:no-arbitrage-and-the-fundamental-theorems:arbitrage}{Arbitrage}
is the stronger and the useful notion: a desk that finds one executes it in
size until the prices move.

\begin{definition}[Box spread]\label{def:dv:no-arbitrage-and-the-fundamental-theorems:box}
A \emph{box spread}\index{box spread} with strikes $K_1<K_2$ and one expiry is
a long call and a short put at $K_1$ together with a short call and a long put
at $K_2$. At expiry it pays $K_2-K_1$ in every state, so its price is
$P(0,T)(K_2-K_1)$ and it is a zero-coupon bond.
\end{definition}

The first two legs are a synthetic long forward at $K_1$, the last two a
synthetic short forward at $K_2$; the difference of two forwards on the same
asset is a sure amount (\cref{fig:dv:no-arbitrage-and-the-fundamental-theorems:box}).
The argument is put--call parity (One Quant Book 1, chapter 25) applied
twice, and like parity it needs European exercise: with American options the
short legs can be assigned early and the box stops being a bond.

\begin{omfigure}
\begin{tikzpicture}
\begin{axis}[width=11cm,height=5cm,
  xlabel={index at expiry},ylabel={payoff (index points)},
  tick label style={font=\small},label style={font=\small},
  xmin=4000,xmax=7000,ymin=-2100,ymax=2100,grid=major,grid style={gray!25},
  xtick={4000,5000,6000,7000},xticklabel style={/pgf/number format/.cd,1000 sep={\,}},
  yticklabel style={/pgf/number format/.cd,1000 sep={\,}},
  legend columns=3,legend style={font=\footnotesize,at={(0.5,-0.3)},anchor=north,draw=none,fill=none,/tikz/every even column/.append style={column sep=8pt}},
  table/col sep=comma]
\addplot[omDef,thick] table[x=s,y=longcall]{figdata/derivatives/01-no-arbitrage-and-the-fundamental-theorems/box_payoff.csv};
\addplot[omDef,thick,dashed] table[x=s,y=shortcall]{figdata/derivatives/01-no-arbitrage-and-the-fundamental-theorems/box_payoff.csv};
\addplot[omProp,thick] table[x=s,y=longput]{figdata/derivatives/01-no-arbitrage-and-the-fundamental-theorems/box_payoff.csv};
\addplot[omProp,thick,dashed] table[x=s,y=shortput]{figdata/derivatives/01-no-arbitrage-and-the-fundamental-theorems/box_payoff.csv};
\addplot[omThm,ultra thick] table[x=s,y=total]{figdata/derivatives/01-no-arbitrage-and-the-fundamental-theorems/box_payoff.csv};
\legend{long 5000 call,short 6000 call,long 6000 put,short 5000 put,the box}
\end{axis}
\end{tikzpicture}

\omcaption{The four legs of a 5\,000--6\,000 \omterm{def:dv:no-arbitrage-and-the-fundamental-theorems:box}{box spread} at expiry. Each leg has
its kink at its strike; the calls' kinks and the puts' kinks cancel, and the sum
is the flat line at 1\,000 points. Data: the chapter's
code.}\label{fig:dv:no-arbitrage-and-the-fundamental-theorems:box}
\end{omfigure}

\begin{example}[Reading the rate off a box]\label{ex:dv:no-arbitrage-and-the-fundamental-theorems:boxrate}
The one-year 5\,000--6\,000 box of the opening trades at 958.90 points. Its
continuously compounded rate is $-\ln(958.90/1000)=4.20\%$. A lender who buys
it pays USD~95\,890 and receives USD~100\,000 in a year (the multiplier is
USD~100 per point); a borrower sells it.
\end{example}

\begin{dated}{2026-09}{The index box-spread market}\label{dat:dv:no-arbitrage-and-the-fundamental-theorems:box}
The exchange that lists the S\&P~500 index options reports USD~929 million of
average daily notional volume in index \omterm{def:dv:no-arbitrage-and-the-fundamental-theorems:box}{box spreads} in 2024. Boxes are mostly
built on 1\,000-point strike distances, a USD~100\,000 notional per box, and the
exchange describes their rates as typically 10 to 55 basis points above
three-month Treasury yields. Every leg is cleared, so the lender's counterparty
is the options clearing house, not the borrower.
\end{dated}

\section{State prices in one period}

\begin{definition}[Arrow--Debreu security and state price]\label{def:dv:no-arbitrage-and-the-fundamental-theorems:ad}
The \emph{Arrow--Debreu security}\index{Arrow--Debreu security} of state $s$
pays 1 if state $s$ occurs and nothing otherwise. Its price $q_s$ is the
\emph{state price}\index{state price} of $s$. A vector $q\in\R^S$ is a
state-price vector for the market if it prices every traded asset:
$p = D^\top q$, that is $p_j=\sum_s q_sD_{sj}$.
\end{definition}

If the Arrow--Debreu securities traded, every claim $g$ would be the portfolio
$\sum_s g_s$ units of security $s$, and would cost $q\cdot g$. They do not trade,
but \omterm{def:dv:no-arbitrage-and-the-fundamental-theorems:ad}{state prices} can still exist: they are the unknowns of the linear system
$D^\top q=p$, one equation per asset.

\begin{theorem}[No arbitrage and positive state prices]\label{thm:dv:no-arbitrage-and-the-fundamental-theorems:ftap1}
In the one-period finite model there is no \omterm{def:dv:no-arbitrage-and-the-fundamental-theorems:arbitrage}{arbitrage} if and only if there is a
state-price vector with $q_s>0$ for every state.
\end{theorem}
\begin{proof}
If $q>0$ prices every asset and $D\theta\ge0$, then $p\cdot\theta=q\cdot
D\theta\ge0$, with equality only if $D\theta=0$: no portfolio is an
\omterm{def:dv:no-arbitrage-and-the-fundamental-theorems:arbitrage}{arbitrage}. Conversely, suppose there is no \omterm{def:dv:no-arbitrage-and-the-fundamental-theorems:arbitrage}{arbitrage}. The set
$C=\{(-p\cdot\theta, D\theta):\theta\in\R^N\}$ is a linear subspace of
$\R^{1+S}$ that meets the non-negative orthant only at zero. By the separating
hyperplane theorem (Farkas' lemma, One Quant Book 4, chapter 23) there is a
vector $(\phi_0,\phi)$ with strictly positive entries orthogonal to $C$:
$-\phi_0\,p\cdot\theta+\phi\cdot D\theta=0$ for all $\theta$. Then
$q=\phi/\phi_0>0$ satisfies $p=D^\top q$.
\end{proof}

\omterm{def:dv:no-arbitrage-and-the-fundamental-theorems:ad}{State prices} carry two familiar objects. If a riskless bond paying 1 in every
state trades, its price is $\sum_sq_s=P(0,T)$, and the normalised
$\pi_s=q_s/P(0,T)$ are positive weights summing to one: a probability, under
which every asset's price is its discounted expected payoff,
$p_j=P(0,T)\sum_s\pi_sD_{sj}$. That probability is the risk-neutral measure of
One Quant Book 4, chapter 5, the equivalent martingale measure with the bank
account as numeraire: prices divided by the bond are martingales. Nothing in it
says that states are likely or unlikely; it is a normalisation of prices.

\begin{definition}[State-price density]\label{def:dv:no-arbitrage-and-the-fundamental-theorems:spd}
Given the real-world probabilities $\P(s)>0$ of the states, the
\emph{state-price density}\index{state-price density} is $m_s=q_s/\P(s)$; every
price is an expectation under $\P$: $p_j=\E[mD_j]$. It is the ratio of the
risk-neutral to the real-world probability times the discount factor, and it is
high in the states in which investors most value a payoff.
\end{definition}

\begin{example}[A binomial market]\label{ex:dv:no-arbitrage-and-the-fundamental-theorems:binomial}
A bond costs 0.98 and pays 1; a share costs 100 and pays 120 or 90. The
equations $q_u+q_d=0.98$ and $120q_u+90q_d=100$ give $q_u=0.3933$, $q_d=0.5867$,
both positive: no \omterm{def:dv:no-arbitrage-and-the-fundamental-theorems:arbitrage}{arbitrage}. The risk-neutral probability of the up state is
$0.3933/0.98=0.4014$, whatever its real-world probability. A call struck at 100
pays $(20,0)$ and is worth $20q_u=7.8667$.
\end{example}

\section{Martingale measures and the first fundamental theorem}

Over several periods, a portfolio is rebalanced at each date $t_k$ according to
what is known then. Let $S_t$ be the vector of risky prices, $B_t$ the bank
account, and $\theta_k$ the holdings chosen at $t_k$ and kept until $t_{k+1}$.

\begin{definition}[Self-financing strategy]\label{def:dv:no-arbitrage-and-the-fundamental-theorems:sf}
A trading strategy $(\theta_k)$, adapted to the information available at each
date, is a \emph{self-financing strategy}\index{self-financing strategy} if
nothing is added or withdrawn between $0$ and $T$: every rebalancing is paid for
by the portfolio itself, so its value changes only through price changes,
$V_{t_{k+1}}-V_{t_k}=\theta_k\cdot(S_{t_{k+1}}-S_{t_k})$ in discounted units. In
continuous time, $dV_t=\theta_t\,dS_t$ with a stochastic integral.
\end{definition}

An \omterm{def:dv:no-arbitrage-and-the-fundamental-theorems:arbitrage}{arbitrage} over several periods is a \omterm{def:dv:no-arbitrage-and-the-fundamental-theorems:sf}{self-financing strategy} with $V_0\le0$,
$V_T\ge0$ and $\P(V_T>0)>0$ (or $V_0<0$). The one-period result extends date
by date.

\begin{theorem}[Fundamental theorem of asset pricing]\label{thm:dv:no-arbitrage-and-the-fundamental-theorems:ftap}
In a market with finitely many trading dates, the \emph{fundamental theorem of
asset pricing}\index{fundamental theorem of asset pricing} states that there is
no \omterm{def:dv:no-arbitrage-and-the-fundamental-theorems:arbitrage}{arbitrage} if and only if there exists an equivalent martingale measure: a
probability $\mathbb Q$ equivalent to $\P$ under which every discounted price
$S_t/B_t$ is a martingale. The arbitrage-free price of a claim $g$ paid at $T$
is then $B_0\E^{\mathbb Q}[g/B_T]$ for such a $\mathbb Q$.
\end{theorem}
\admitted

For a finite state space the proof is the one-period theorem applied on every
node of the information tree, the \omterm{def:dv:no-arbitrage-and-the-fundamental-theorems:ad}{state prices} of each node giving the
conditional probabilities of $\mathbb Q$ (Harrison and Kreps, then Harrison and
Pliska). For general discrete time it is due to Dalang, Morton and Willinger;
in continuous time ``no \omterm{def:dv:no-arbitrage-and-the-fundamental-theorems:arbitrage}{arbitrage}'' must be strengthened to ``no free lunch
with vanishing risk'' and the martingale property weakened to a sigma-martingale
(Delbaen and Schachermayer); the measure-theoretic statements and Girsanov's
theorem, which builds $\mathbb Q$ in diffusion models, are in One Quant Book 4,
chapters 1 and 5. What the desk needs is the direction of use: to price, find
a measure under which discounted traded prices are martingales, and take the
expectation.

\begin{remark}[Why the bank account]\label{rem:dv:no-arbitrage-and-the-fundamental-theorems:numeraire}
Nothing forces the bank account as unit. Any traded asset with positive price
can serve as numeraire $\mathcal N_t$: under its measure $\mathbb Q^{\mathcal N}$
prices divided by $\mathcal N_t$ are martingales. Choosing the numeraire well
(the zero-coupon bond for a payoff at $T$, the asset itself for a payoff in
shares) removes a stochastic discount from an expectation; chapter 3 uses it to
derive the Black--Scholes formula in two lines.
\end{remark}

\section{Completeness, replication and the second theorem}

\begin{definition}[Complete market]\label{def:dv:no-arbitrage-and-the-fundamental-theorems:complete}
A \emph{complete market}\index{complete market} is one in which every
\omterm{def:dv:no-arbitrage-and-the-fundamental-theorems:claim}{contingent claim} has a replicating \omterm{def:dv:no-arbitrage-and-the-fundamental-theorems:sf}{self-financing strategy}. In the one-period
finite model, the market is complete if and only if $D$ has rank $S$.
\end{definition}

\begin{theorem}[Second fundamental theorem]\label{thm:dv:no-arbitrage-and-the-fundamental-theorems:ftap2}
In an arbitrage-free market with finitely many dates and states, the market is
complete if and only if the equivalent martingale measure is unique.
\end{theorem}
\begin{proof}[Partial proof (one period)]
The state-price vectors are the solutions of $D^\top q=p$ with $q>0$. If $D$
has rank $S$, $D^\top$ is injective and the solution is unique. If the rank is
less than $S$, the kernel of $D^\top$ contains some $z\ne0$, and $q+\varepsilon z$
remains a positive solution for small $\varepsilon$: the measure is not unique.
\end{proof}

A \omterm{def:dv:no-arbitrage-and-the-fundamental-theorems:complete}{complete market} prices every claim: the price is the cost of the \omterm{def:dv:no-arbitrage-and-the-fundamental-theorems:claim}{replicating
portfolio}, and holding that portfolio against a sale of the claim leaves no
risk. This is what makes option pricing a trade and not an opinion. The
binomial market of \cref{ex:dv:no-arbitrage-and-the-fundamental-theorems:binomial}
is complete; chapter 2 builds a whole pricing theory from it. Add a third state
and it is not.

\begin{omfigure}
\begin{tikzpicture}[font=\small,
  nd/.style={draw=omDef,thick,rounded corners=2pt,fill=omDef!6,align=center,minimum width=2.4cm,inner sep=3pt}]
\node[nd] (o) at (0,0) {today\\bond 0.98, share 100};
\node[nd] (u) at (6,1.7) {state $u$: share 120};
\node[nd] (m) at (6,0) {state $m$: share 100};
\node[nd] (d) at (6,-1.7) {state $d$: share 80};
\draw[-{Stealth},thick,gray] (o.east) -- (u.west);
\draw[-{Stealth},thick,gray] (o.east) -- (m.west);
\draw[-{Stealth},thick,gray] (o.east) -- (d.west);
\node[anchor=west,align=left,font=\footnotesize] at (8.3,1.7) {$q_u\in(0.10,\,0.54)$};
\node[anchor=west,align=left,font=\footnotesize] at (8.3,0) {$q_m\in(0,\,0.88)$};
\node[anchor=west,align=left,font=\footnotesize] at (8.3,-1.7) {$q_d\in(0,\,0.44)$};
\node[align=center,font=\footnotesize,text=omThm] at (3,-2.6) {vertices $(0.10,\,0.88,\,0)$ and $(0.54,\,0,\,0.44)$:\\every positive state-price vector lies strictly between them};
\end{tikzpicture}

\omcaption{The trinomial market: two assets, three states. The two equations
$q_u+q_m+q_d=0.98$ and $120q_u+100q_m+80q_d=100$ leave one degree of freedom, a
segment of state-price vectors whose end points each give one state a zero
price. Illustrative numbers.}\label{fig:dv:no-arbitrage-and-the-fundamental-theorems:tree}
\end{omfigure}

\section{Price bounds when the market is incomplete}

In the trinomial market of \cref{fig:dv:no-arbitrage-and-the-fundamental-theorems:tree}
a call struck at 100 pays $(20,0,0)$. Its price must be $20q_u$ for some
positive state-price vector, so anything strictly between $20\times0.10=2$ and
$20\times0.54=10.80$ is consistent with no \omterm{def:dv:no-arbitrage-and-the-fundamental-theorems:arbitrage}{arbitrage}. The market has not chosen.

\begin{definition}[Super-replication price]\label{def:dv:no-arbitrage-and-the-fundamental-theorems:super}
The \emph{super-replication price}\index{super-replication price} of a claim
$g$ is the lowest cost of a portfolio of traded assets that pays at least $g$ in
every state: $\inf\{p\cdot\theta: D\theta\ge g\}$. The sub-replication price is the
highest cost of one that pays at most $g$.
\end{definition}

\begin{proposition}[The no-arbitrage interval]\label{prop:dv:no-arbitrage-and-the-fundamental-theorems:interval}
In an arbitrage-free one-period market, the \omterm{def:dv:no-arbitrage-and-the-fundamental-theorems:super}{super-replication price} of $g$ equals
$\max q\cdot g$ over state-price vectors $q\ge0$ with $D^\top q=p$, and the
sub-replication price equals the minimum. A new asset paying $g$ can be added at
price $c$ without creating an \omterm{def:dv:no-arbitrage-and-the-fundamental-theorems:arbitrage}{arbitrage} if and only if $c$ lies strictly between
them, or $g$ is replicable and $c$ is its price.
\end{proposition}
\begin{proof}
The first statement is linear-programming duality: the dual of $\min\{p\cdot
\theta:D\theta\ge g\}$ is $\max\{q\cdot g:D^\top q=p,\ q\ge0\}$, and both optima
are attained and equal when the feasible sets are non-empty. For the second,
apply \cref{thm:dv:no-arbitrage-and-the-fundamental-theorems:ftap1} to the
enlarged market: a positive state-price vector exists for it exactly when
$c=q\cdot g$ for some $q>0$ pricing the old assets, and the values $q\cdot g$ over
the open set of positive $q$ fill the open interval (or a single point when $g$
is replicable).
\end{proof}

At the end of the interval the claim is still an \omterm{def:dv:no-arbitrage-and-the-fundamental-theorems:arbitrage}{arbitrage}: at 10.80 the call
can be sold and super-replicated at zero net cost by $0.5$ share and $-40$
bonds, and the hedged position pays $(0,10,0)$, nothing in two states and 10 in
the third. \Cref{fig:dv:no-arbitrage-and-the-fundamental-theorems:bounds} shows
the interval by strike: wide at the money, where the call's payoff depends on
the state the two assets cannot tell apart, closed at the ends of the strike
range, where the call is replicable.

\begin{omfigure}
\begin{tikzpicture}
\begin{axis}[width=10.5cm,height=5cm,
  xlabel={strike},ylabel={call price},
  tick label style={font=\small},label style={font=\small},
  xmin=80,xmax=120,ymin=0,ymax=23,grid=major,grid style={gray!25},
  legend columns=3,legend style={font=\footnotesize,at={(0.5,-0.3)},anchor=north,draw=none,fill=none,/tikz/every even column/.append style={column sep=8pt}},
  table/col sep=comma]
\addplot[name path=up,omThm,thick] table[x=strike,y=upper]{figdata/derivatives/01-no-arbitrage-and-the-fundamental-theorems/bounds.csv};
\addplot[name path=lo,omDef,thick] table[x=strike,y=lower]{figdata/derivatives/01-no-arbitrage-and-the-fundamental-theorems/bounds.csv};
\addplot[omProp!25] fill between[of=up and lo];
\addplot[omProp,thick,dashed] table[x=strike,y=complete]{figdata/derivatives/01-no-arbitrage-and-the-fundamental-theorems/bounds.csv};
\legend{super-replication,sub-replication,,once the 100 call trades at 6}
\end{axis}
\end{tikzpicture}

\omcaption{The \omterm{def:dv:no-arbitrage-and-the-fundamental-theorems:arbitrage}{no-arbitrage} interval of a call in the trinomial market, by
strike. Once one option trades (the 100 call at 6.00) the market is complete and
every other strike has a single price, the dashed line, inside the band. Data:
the tutorial.}\label{fig:dv:no-arbitrage-and-the-fundamental-theorems:bounds}
\end{omfigure}

\begin{method}[Static bounds on a quoted chain]\label{met:dv:no-arbitrage-and-the-fundamental-theorems:chain}
For European calls on one underlying with forward $F$ and discount factor $P$,
without any model:
\begin{enumerate}
\item $P\max(F-K,0)\le C(K)\le PF$; for puts $P\max(K-F,0)\le P(K)\le PK$;
\item $C$ decreases in $K$ and $C(K_1)-C(K_2)\le P(K_2-K_1)$ (the call spread
pays at most the strike distance);
\item $C$ is convex in $K$: a butterfly costs at least zero;
\item $C(K)-P(K)=P(F-K)$ at every strike, and every box costs $P(K_2-K_1)$.
\end{enumerate}
With bid and ask quotes each check is made at executable prices; a failure is an
\omterm{def:dv:no-arbitrage-and-the-fundamental-theorems:arbitrage}{arbitrage} whose portfolio the check names
(\cref{fig:dv:no-arbitrage-and-the-fundamental-theorems:region}).
\end{method}

\begin{omfigure}
\begin{tikzpicture}
\begin{axis}[width=10.5cm,height=5cm,
  xlabel={strike},ylabel={call price (index points)},
  tick label style={font=\small},label style={font=\small},
  xmin=3000,xmax=8000,ymin=0,ymax=6200,grid=major,grid style={gray!25},xtick={3000,4000,5000,6000,7000,8000},
  xticklabel style={/pgf/number format/.cd,1000 sep={\,}},yticklabel style={/pgf/number format/.cd,1000 sep={\,}},
  legend columns=3,legend style={font=\footnotesize,at={(0.5,-0.3)},anchor=north,draw=none,fill=none,/tikz/every even column/.append style={column sep=8pt}},
  table/col sep=comma]
\addplot[omThm,thick,dashed] table[x=strike,y=upper]{figdata/derivatives/01-no-arbitrage-and-the-fundamental-theorems/call_region.csv};
\addplot[omDef,thick,dashed] table[x=strike,y=lower]{figdata/derivatives/01-no-arbitrage-and-the-fundamental-theorems/call_region.csv};
\addplot[omProp,very thick] table[x=strike,y=call]{figdata/derivatives/01-no-arbitrage-and-the-fundamental-theorems/call_region.csv};
\legend{discounted forward,discounted intrinsic value,a smile-consistent call curve}
\end{axis}
\end{tikzpicture}

\omcaption{Where a call price can lie: between the discounted intrinsic value
$P\max(F-K,0)$ and the discounted forward $PF$, decreasing and convex in the
strike. The solid curve is the chain of the weekend problem (forward 5\,850,
one year, a skewed smile). Illustrative
parameters.}\label{fig:dv:no-arbitrage-and-the-fundamental-theorems:region}
\end{omfigure}

\section{Tutorial: state prices, bounds and a box-spread rate}\label{tut:dv:no-arbitrage-and-the-fundamental-theorems:tut}

\begin{tutorial}
\textbf{Goal.} Solve for \omterm{def:dv:no-arbitrage-and-the-fundamental-theorems:ad}{state prices}, compute the \omterm{def:dv:no-arbitrage-and-the-fundamental-theorems:arbitrage}{no-arbitrage} interval of a
claim in an incomplete market, find an \omterm{def:dv:no-arbitrage-and-the-fundamental-theorems:arbitrage}{arbitrage} portfolio when a price leaves
the interval, and read a rate off a quoted box. \textbf{End state:}
\cref{fig:dv:no-arbitrage-and-the-fundamental-theorems:bounds} and the numbers of
\cref{ex:dv:no-arbitrage-and-the-fundamental-theorems:binomial} and of the
weekend problem.
\begin{tutsteps}
\item \textbf{Vertices of the state-price set.} Every vertex solves the system on
a support of as many states as there are assets; the bounds of a claim are
attained at vertices.
\omcode{code/derivatives/01-no-arbitrage-and-the-fundamental-theorems/python/dv_arbitrage.py}{33}{52}{State-price vertices and the no-arbitrage bounds of a claim.}\label{lst:dv:no-arbitrage-and-the-fundamental-theorems:vertices}
\item \textbf{The primal side.} The cheapest super-replicating portfolio, and a
search for an \omterm{def:dv:no-arbitrage-and-the-fundamental-theorems:arbitrage}{arbitrage} portfolio, are linear programmes; the build's
\texttt{lp\_max} solves them exactly by enumerating basic solutions.
\omcode{code/derivatives/01-no-arbitrage-and-the-fundamental-theorems/python/dv_arbitrage.py}{55}{74}{Super-replication and an arbitrage search as linear programmes.}\label{lst:dv:no-arbitrage-and-the-fundamental-theorems:lp}
\item \textbf{Run} the tests: the upper bound equals the super-replication cost
for every strike, the call spread 90--110 has a one-point interval, and pricing
the 100 call at 11.00 returns the \omterm{def:dv:no-arbitrage-and-the-fundamental-theorems:arbitrage}{arbitrage} $(-40, 0.5, -1)$ with an edge of 0.20.
\item \textbf{The box.} Price the four legs at executable prices and invert:
\omcode{code/firm/arbcheck/firm_arbcheck.py}{92}{102}{The cost of buying or selling a box, and the rate it implies.}\label{lst:dv:no-arbitrage-and-the-fundamental-theorems:box}
\end{tutsteps}
\textbf{What to change next.} Add a fourth state and a second option and watch
the interval of a third option shrink; then quote the four box legs one tick
wider and see how far the lending and borrowing rates move apart.
\end{tutorial}

\section{Build: the static-arbitrage checker}\label{bld:dv:no-arbitrage-and-the-fundamental-theorems:arbcheck}

\begin{build}
\bfield{Purpose} Before any model touches a chain of quotes, the miniature firm
checks that the chain itself is free of static \omterm{def:dv:no-arbitrage-and-the-fundamental-theorems:arbitrage}{arbitrage}, and reports the
portfolio and the edge of every violation. The same checks guard the surface
fits of chapters 7 and 8.
\bfield{Interface} \texttt{Quote(strike, right, bid, ask)};
\texttt{check\_chain(quotes, forward, df) -> list[Violation]} with
\texttt{Violation(kind, strikes, edge, legs)}; \texttt{box\_cost},
\texttt{box\_rate}, \texttt{check\_box}; \texttt{lp\_max(c, a\_ub, b\_ub)}.
\bfield{Rules} Buy at the ask, sell at the bid; an edge is cash received today by a
portfolio whose payoff is never negative; convexity is checked on unequal
strike spacing with the weights $(K_3-K_2)/(K_3-K_1)$ and $(K_2-K_1)/(K_3-K_1)$.
\bfield{Acceptance tests} \texttt{code/firm/arbcheck/tests/}: a Black--Scholes
chain passes; a mispriced strike is caught by the convexity check with a
positive edge; an option below its discounted intrinsic value is caught with the
right edge; the box rate of a clean chain equals the chain's rate.
\bfield{Stretch} Replace the rule list by one linear programme over all
static portfolios of the chain (options, forward, bond), which also finds
violations that combine more than three strikes.
\end{build}

\begin{omsources}
\item J.~M.~Harrison and D.~M.~Kreps, ``Martingales and arbitrage in multiperiod
securities markets'', \emph{Journal of Economic Theory} 20 (1979) 381--408.
\item J.~M.~Harrison and S.~R.~Pliska, ``Martingales and stochastic integrals in
the theory of continuous trading'', \emph{Stochastic Processes and their
Applications} 11 (1981) 215--260.
\item F.~Delbaen and W.~Schachermayer, ``A general version of the fundamental
theorem of asset pricing'', \emph{Mathematische Annalen} 300 (1994) 463--520.
\item K.~J.~Arrow, ``The role of securities in the optimal allocation of
risk-bearing'', \emph{Review of Economic Studies} 31 (1964) 91--96.
\item J.~H.~van Binsbergen, W.~F.~Diamond and M.~Grotteria, ``Risk-free
interest rates'', \emph{Journal of Financial Economics} 143 (2022) 1--29.
\item Cboe, SPX box spreads (product material, 2024--2025).
\end{omsources}

\section{Exercises}

\begin{exercise}[$\star$]\label{exo:dv:no-arbitrage-and-the-fundamental-theorems:1}
A bond costs 0.99 and pays 1; a share costs 100 and pays 115 or 95. Find the \omterm{def:dv:no-arbitrage-and-the-fundamental-theorems:ad}{state
prices}, the risk-neutral probability of the up state, and the price of a call
struck at 100.
\end{exercise}

\begin{exercise}[$\star$]\label{exo:dv:no-arbitrage-and-the-fundamental-theorems:2}
Check, for an index at expiry below 5\,000, between the strikes, and above
6\,000, that the 5\,000--6\,000 box pays 1\,000 points.
\end{exercise}

\begin{exercise}[$\star$]\label{exo:dv:no-arbitrage-and-the-fundamental-theorems:3}
In the trinomial market, a put struck at 100 is offered at 9.00. Is there an
\omterm{def:dv:no-arbitrage-and-the-fundamental-theorems:arbitrage}{arbitrage}? If so, give the portfolio and its edge.
\end{exercise}

\begin{exercise}[$\star\star$]\label{exo:dv:no-arbitrage-and-the-fundamental-theorems:4}
Build a two-state market that satisfies the \omterm{def:dv:no-arbitrage-and-the-fundamental-theorems:loop}{law of one price} and has an
\omterm{def:dv:no-arbitrage-and-the-fundamental-theorems:arbitrage}{arbitrage}. Which of the conditions of
\cref{thm:dv:no-arbitrage-and-the-fundamental-theorems:ftap1} fails?
\end{exercise}

\begin{exercise}[$\star\star$]\label{exo:dv:no-arbitrage-and-the-fundamental-theorems:5}
In the trinomial market completed by the 100 call at 6.00, the \omterm{def:dv:no-arbitrage-and-the-fundamental-theorems:ad}{state prices} are
$(0.30, 0.48, 0.20)$. With real-world probabilities $(0.45, 0.45, 0.10)$, give
the risk-neutral probabilities and the \omterm{def:dv:no-arbitrage-and-the-fundamental-theorems:spd}{state-price density}, and say what the
density's shape means.
\end{exercise}

\begin{exercise}[$\star\star$]\label{exo:dv:no-arbitrage-and-the-fundamental-theorems:6}
Give the \omterm{def:dv:no-arbitrage-and-the-fundamental-theorems:arbitrage}{no-arbitrage} interval of a digital that pays 1 in the up state of the
trinomial market, and the super-replicating portfolio at its upper end.
\end{exercise}

\begin{exercise}[$\star\star\star$]\label{exo:dv:no-arbitrage-and-the-fundamental-theorems:7}
\emph{Coding.} With \texttt{price\_bounds}, compute the interval of the 90--110
call spread in the trinomial market. Explain the result.
\end{exercise}

\begin{exercise}[$\star\star\star$]\label{exo:dv:no-arbitrage-and-the-fundamental-theorems:8}
\emph{Find the flaw.} ``In the trinomial market the 100 call can trade anywhere
between 2 and 10.80. My model says it is worth 6, so I buy it at 5 and hedge: that
is a riskless profit of 1.'' Correct the reasoning.
\end{exercise}

\section{Problem: The Box-Spread Rate}

\begin{problem}[{Weekend problem --- lending cash through four options}]\label{pb:dv:no-arbitrage-and-the-fundamental-theorems:1}
A treasurer wants to lend USD~10 million for one year through index \omterm{def:dv:no-arbitrage-and-the-fundamental-theorems:box}{box spreads}
and compares the quotes with the overnight rate, which she expects to compound
to 4.00\% (continuously compounded) over the year. The one-year options on the
index, European and cash-settled with a USD~100 multiplier, are quoted (index
points): 5\,000 call 971.80--973.40, 6\,000 call 324.60--326.20, 5\,000 put
156.70--158.30, 6\,000 put 468.40--470.00. The quotes are illustrative.

\medskip\noindent\textbf{Part I --- The four legs.}
\begin{enumerate}
\item Write the payoff of each leg at expiry and show that the total is 1\,000
points in every state.
\item What does one box pay in dollars?
\item What does the treasurer pay for one box bought leg by leg at the quotes?
\item What would a borrower receive for one box sold leg by leg?
\item Which risks does the box remove, and which remain?
\end{enumerate}

\medskip\noindent\textbf{Part II --- Rates.}
\begin{enumerate}[resume]
\item Give the lending rate implied by buying at the quotes.
\item Give the borrowing rate implied by selling at the quotes.
\item Give the mid price of the box and its rate.
\item How many dollars of interest does the treasurer earn per box bought leg by leg?
\item Give the spreads of the lending and borrowing rates over the overnight rate, in basis points.
\end{enumerate}

\medskip\noindent\textbf{Part III --- Packages and \omterm{def:dv:no-arbitrage-and-the-fundamental-theorems:arbitrage}{arbitrage}.}
\begin{enumerate}[resume]
\item Why do boxes trade as a package at one price rather than leg by leg?
\item Using the overnight rate as the discount rate, is there an \omterm{def:dv:no-arbitrage-and-the-fundamental-theorems:arbitrage}{arbitrage} in the quotes?
\item At what package price would the treasurer earn exactly the overnight rate?
\item What would change if the options were American, as single-stock options are?
\item Give one reason why box rates can exceed Treasury bill yields of the same maturity.
\end{enumerate}

\medskip\noindent\textbf{Part IV --- Judgement.}
\begin{enumerate}[resume]
\item Who borrows by selling boxes, and why?
\item How many boxes does the treasurer need, and what does she pay at the mid price?
\item What does she hold if the clearing member she uses defaults?
\item State the \emph{named result}: the box's mid implied rate and its spread over the overnight rate.
\item In one sentence: why is a \omterm{def:dv:no-arbitrage-and-the-fundamental-theorems:box}{box spread} a zero-coupon bond?
\end{enumerate}
\end{problem}

\section{Interview questions}

\begin{interviewq}[$\star$ \iqroles{trader, researcher}]\label{iq:dv:no-arbitrage-and-the-fundamental-theorems:1}
What is a \omterm{def:dv:no-arbitrage-and-the-fundamental-theorems:box}{box spread}, what does it pay, and what does its price tell you?
\end{interviewq}

\begin{interviewq}[$\star$ \iqroles{researcher, bank}]\label{iq:dv:no-arbitrage-and-the-fundamental-theorems:2}
State the \omterm{thm:dv:no-arbitrage-and-the-fundamental-theorems:ftap}{fundamental theorem of asset pricing}, and say what the risk-neutral
probabilities are and are not.
\end{interviewq}

\begin{interviewq}[$\star\star$ \iqroles{trader}]\label{iq:dv:no-arbitrage-and-the-fundamental-theorems:3}
A share trades at 100 and will be at 110 or 95 in a year; the one-year rate is
zero. Price a digital that pays 1 if the share ends at 110.
\end{interviewq}

\begin{interviewq}[$\star\star$ \iqroles{researcher}]\label{iq:dv:no-arbitrage-and-the-fundamental-theorems:4}
Why is a \omterm{def:dv:no-arbitrage-and-the-fundamental-theorems:complete}{complete market} equivalent to a unique martingale measure? Give an
example of an incomplete market and say what happens to option prices in it.
\end{interviewq}

\begin{interviewq}[$\star\star$ \iqroles{trader, risk}]\label{iq:dv:no-arbitrage-and-the-fundamental-theorems:5}
Calls at strikes 90, 100 and 110 have mid prices 12.00, 6.50 and 0.90, each
quoted 0.10 wide. Is there an \omterm{def:dv:no-arbitrage-and-the-fundamental-theorems:arbitrage}{arbitrage}? What would you trade?
\end{interviewq}

\begin{interviewq}[$\star\star\star$ \iqroles{developer}]\label{iq:dv:no-arbitrage-and-the-fundamental-theorems:6}
Design a service that checks 3\,000 listed option quotes on one underlying for
static \omterm{def:dv:no-arbitrage-and-the-fundamental-theorems:arbitrage}{arbitrage} every time a quote changes.
\end{interviewq}
